\section{Resultados e validação}
\subsection{Estabilidade absoluta e relativa}
Foram aprovadas \ChecksPass{} de \ChecksTotal{} verificações
computacionais nos experimentos. A Tabela~\ref{tab:margins}
apresenta os pontos nominais. Routh também classificou os quatro
casos nominais como estáveis, usando Padé $[3/3]$ para D.
A comparação incluiu configurações abaixo, sobre e acima das
fronteiras. Em A, o ganho $K=10$ levou um par ao eixo. Em B,
$K=2$ produziu um polo na origem. Em C, o par imaginário
surgiu em $K=3$. O cálculo independente do lugar das raízes
concordou com esses três pontos de ganho. As
Figs.~\ref{fig:nyquist}--\ref{fig:locus} mostram os diagramas
nos quatro sistemas. A Fig.~\ref{fig:margins} compara as margens
nominais na convenção da tabela.

\begin{table}[!t]
\centering
\caption{Estabilidade e margem nos pontos nominais.}
\label{tab:margins}
\begin{tabular}{cccc}
\toprule
Caso & $(P,N,Z)$ & MG (razão) & MF ($^\circ$)\\
\midrule
\NominalRows
\bottomrule
\end{tabular}
\end{table}

No caso A, a margem de fase da tabela decorre de um cruzamento
em DC. O comparador Python-control reporta a margem de ganho
de 10, mas sua rotina de margens usuais omite esse cruzamento
unitário em DC. A diferença foi registrada nas convenções
experimentais. Em B, a margem de ganho de $0{,}5$ identifica a
fronteira ao reduzir o ganho nominal de 4 para 2. A contagem
$P=1$ e $N=-1$ demonstra a estabilidade em malha fechada.
Em C, inexiste um cruzamento de ganho unitário na frequência
real e a MG é 3. A resposta inicial inversa produz um mínimo
amostrado de $-0{,}19487$, embora a amostra final se aproxime
de $1/3$ no intervalo avaliado. A Fig.~\ref{fig:step} mostra a
resposta ao degrau dos quatro casos.

O erro relativo máximo da resposta complexa nominal frente à
referência foi de $6{,}22\times10^{-16}$. O maior erro absoluto
nos conjuntos de raízes por ganho foi de $3{,}56\times10^{-14}$
e a maior diferença no degrau foi de $7{,}25\times10^{-14}$
nas quatro plantas, dentro das tolerâncias do protocolo.
Para D, a referência temporal representa o mesmo modelo racional com Padé $[3/3]$. O acordo nesse degrau mede o algoritmo de resposta
temporal e não mede erro físico do Padé.

\begin{figure*}[!t]
\centering
\figfile{0.95\textwidth}{nyquist-abcd.pdf}
\caption{Contornos de Nyquist das plantas A--D nos pontos nominais.
A cruz assinala $-1$. O caso B exige um envoltório anti-horário.}
\label{fig:nyquist}
\end{figure*}

\begin{figure*}[!t]
\centering
\figfile{0.97\textwidth}{bode-abcd.pdf}
\caption{Magnitude e fase de A--D. O caso D usa o atraso exponencial.
A faixa exibida de frequências é finita.}
\label{fig:bode}
\end{figure*}

\begin{figure*}[!t]
\centering
\figfile{0.95\textwidth}{root-locus-abcd.pdf}
\caption{Lugar das raízes em A--D. Em D, os ramos são os
do modelo racional de Padé $[3/3]$.}
\label{fig:locus}
\end{figure*}

\begin{figure*}[!t]
\centering
\figfile{0.96\textwidth}{stability-margins.pdf}
\caption{Margens dos quatro sistemas nominais. Para D,
selecionou-se o primeiro cruzamento de fase positivo em frequência.}
\label{fig:margins}
\end{figure*}

\begin{figure}[!t]
\centering
\figfile{\columnwidth}{step-response.pdf}
\caption{Degrau unitário de malha fechada. A curva C apresenta
deslocamento inicial negativo. Para D usa-se Padé $[3/3]$.}
\label{fig:step}
\end{figure}

\subsection{Atraso exato e Padé}
Em D, o menor cruzamento de módulo unitário positivo fica em
$\sqrt3$ rad/s. O cálculo encontrou a fronteira em
\DelayCritical{} s, em concordância com \eqref{eq:delay}.
No ponto nominal, a MF é $40{,}61^\circ$. O primeiro cruzamento
de fase positivo ocorre em $2{,}44822$ rad/s e fornece a MG
linear de $1{,}32229$. Em frequências mais altas existem
cruzamentos adicionais, com margens correspondentes ao módulo
cada vez menor da planta.

Em $\tau=5$ s, a primeira versão com refinamento global
encontrou quatro envoltórios, mas a variação de argumento
entre amostras excedeu seu critério e o resultado ficou
indeterminado. A bissecção localizada resolveu a transição
com a configuração padrão. Contaram-se $P=0$, $N=4$ e
$Z=4$, confirmados por dois pares à direita na solução por
Lambert $W$. A execução arquivou o diagnóstico original e
os valores antes e depois do refinamento.

A Tabela~\ref{tab:pade} compara a primeira fronteira dos modelos
racionais com o atraso exato. Em $1{,}01\tau_*$, as ordens
1 e 2 ainda indicam estabilidade, embora o contorno exato
conclua instabilidade. As ordens 3 e 5 recuperam a classificação
nesse ponto. O erro complexo e a discrepância de classificação
podem ser observados na Fig.~\ref{fig:pade}. A ordem adequada
depende da banda e da distância da fronteira.

\begin{table}[!t]
\centering
\caption{Atraso crítico estimado por Padé para D.}
\label{tab:pade}
\begin{tabular}{cccc}
\toprule
Ordem & $\hat\tau_*$ (s) & Erro relativo & Em $1{,}01\tau_*$\\
\midrule
\DelayRows
\bottomrule
\end{tabular}
\end{table}

\begin{figure*}[!tp]
\centering
\figfile{0.88\textwidth}{delay-pade.pdf}
\caption{Diferenças de frequência e de atraso crítico no caso D.
A referência usa a exponencial.}
\label{fig:pade}
\end{figure*}

\subsection{Desempenho e erro numérico}
As medições foram executadas em macOS 26.6.2, em um Apple M5 com
10 núcleos lógicos e 16 GiB. Usaram-se dois aquecimentos e
11 repetições por ordem e método, num total de 132 medições.
A varredura do lugar das raízes contém 48 ganhos e a de Bode,
1001 frequências. A semente da perturbação dos coeficientes
é 24112026. As amostras, os valores de CPU e parede, a média,
a mediana e o desvio-padrão estão preservados para auditoria.

\begin{table}[!t]
\centering
\caption{Mediana do tempo de CPU em milissegundos.}
\label{tab:cpu}
\begin{tabular}{lrrr}
\toprule
Método & $n=5$ & $n=10$ & $n=20$\\
\midrule
\CpuRows
\bottomrule
\end{tabular}
\end{table}

A Tabela~\ref{tab:cpu} e a Fig.~\ref{fig:bench} registram o crescimento
observado no protocolo. A carga de Routh é uma tabela por
polinômio. Lugar das raízes calcula conjuntos de raízes e
associa trajetórias para 48 ganhos. Nyquist avalia múltiplos
contornos e conta envoltórios. Bode calcula margens e amostra
duas curvas. Essas diferenças de tarefa condicionam a interpretação
das medições. Os valores de memória usam \texttt{tracemalloc}
em execução separada, refletindo o pico de alocações rastreadas
pelo Python.

\begin{figure*}[!tp]
\centering
\figfile{0.88\textwidth}{scalability.pdf}
\caption{Escalabilidade no protocolo de 11 repetições por combinação.
A dispersão foi medida com gráficos construídos dos dados brutos.}
\label{fig:bench}
\end{figure*}

O experimento de precisão comparou duas famílias de raízes
construídas, uma dispersa e uma agrupada. Em ordem 20, os erros
absolutos máximos em relação às raízes de projeto foram
$0{,}13280$ e $0{,}31928$ s$^{-1}$. Os máximos de
condicionamento relativo foram $7{,}59\times10^{14}$ e
$2{,}27\times10^{24}$. As duas famílias exibiram resíduos
escalados inferiores a $3\times10^{-16}$. Esses valores mostram
que uma raiz pode satisfazer numericamente o polinômio com
resíduo pequeno e ainda estar distante da raiz de projeto.
A referência em precisão estendida separa o arredondamento
dos coeficientes do erro do solucionador no polinômio
armazenado. A experiência repetiu o cálculo após perturbação
relativa de $10^{-10}$ nos coeficientes.
Na Fig.~\ref{fig:precision}, a dispersão dos resultados depende
da distribuição das raízes e da ordem considerada.

\begin{figure*}[!tp]
\centering
\figfile{0.88\textwidth}{numerical-precision.pdf}
\caption{Erro de localização e condicionamento para as famílias
dispersa e agrupada em três ordens.}
\label{fig:precision}
\end{figure*}
